\chapter{DistributionFitter}
\label{chap:distribution_fitter}

\section{Purpose}

The \texttt{DistributionFitter} class is used to fit one or more candidate
probability distributions to an observed dataset and to compare them using
goodness-of-fit metrics. In actuarial practice, this step typically
precedes simulation: the analyst inspects historical claim counts or
severities, selects a plausible distributional form, and validates the
choice using diagnostic statistics.

\section{Supported Distributions}

\texttt{DistributionFitter} supports the same distributions exposed by
\texttt{actstats}, summarised below.

\begin{center}
\begin{tabular}{ll}
\toprule
\textbf{Type} & \textbf{Distributions} \\
\midrule
Continuous severity & \texttt{normal}, \texttt{logistic},
                      \texttt{exponential}, \texttt{gamma}, \texttt{beta},
                      \texttt{lognormal}, \texttt{weibull},
                      \texttt{pareto}, \texttt{uniform} \\
Discrete frequency  & \texttt{poisson}, \texttt{negative binomial} \\
\bottomrule
\end{tabular}
\end{center}

The distributions used in any particular fitting run can be controlled
through the \texttt{distributions} argument or via a YAML configuration file.

\section{Goodness-of-Fit Metrics}

For each fitted distribution, \texttt{DistributionFitter} computes:

\begin{itemize}[leftmargin=*]
    \item Log-likelihood
    \item Akaike Information Criterion (AIC)
    \item Bayesian Information Criterion (BIC)
    \item Chi-square statistic
    \item Kolmogorov--Smirnov statistic
\end{itemize}

The information criteria are defined as
\[
    \mathrm{AIC} = 2k - 2\ln L, \qquad
    \mathrm{BIC} = k \ln n - 2 \ln L,
\]
where $k$ is the number of estimated parameters, $n$ is the sample size,
and $L$ is the likelihood evaluated at the maximum-likelihood estimate.

\subsection{Pearson Chi-Square Statistic}

\subsubsection{Definition}
Partition the real line into $k$ bins $(a_0, a_1], (a_1, a_2], \dots, (a_{k-1}, a_k]$ with
$a_0 = -\infty$, $a_k = +\infty$. Let
\[
O_i = \#\{\, x_j : a_{i-1} < x_j \le a_i \,\}, \qquad
E_i = n \Big[F(a_i;\hat\theta) - F(a_{i-1};\hat\theta)\Big]
\]
be the observed and expected counts in bin $i$, where $n$ is the sample size. The Pearson
chi-square statistic is
\[
\boxed{\;\chi^2 = \sum_{i=1}^{k} \frac{(O_i - E_i)^2}{E_i}\;}
\]
Under the null hypothesis that the data were generated by $F(\cdot;\hat\theta)$, $\chi^2$ is
asymptotically distributed as $\chi^2_{k-1-m}$, where $m$ is the number of parameters
estimated from the data (a further $-1$ degree of freedom is lost for each such parameter,
beyond the $k-1$ already lost to the constraint $\sum_i O_i = \sum_i E_i = n$).

\subsubsection{Bin construction: equiprobable bins from the fitted quantile function}
The choice of bin edges $\{a_i\}$ is not unique, and a poor choice can produce misleading
results. Two approaches are common:
\begin{itemize}[leftmargin=1.4em]
  \item \textbf{Data-driven bins} (e.g.\ \texttt{numpy.histogram\_bin\_edges}): edges chosen
  from the empirical range of the sample. This makes bin \emph{expected} counts $E_i$
  data-dependent and can produce empty or near-empty expected bins in the tails, where the
  $\chi^2$ approximation breaks down ($E_i \to 0$ makes each term $(O_i-E_i)^2/E_i$
  ill-behaved).
  \item \textbf{Equiprobable bins} (used here): edges are chosen so that \emph{under the
  fitted distribution}, each bin has equal probability $1/k$. This is the standard
  construction recommended in the goodness-of-fit literature because it guarantees
  $E_i = n/k$ exactly for every bin (up to floating point), for any fitted distribution.
\end{itemize}
Equiprobable interior edges are obtained from the fitted quantile function (inverse CDF):
\[
a_i = F^{-1}\!\left(\frac{i}{k};\hat\theta\right), \qquad i = 1, \dots, k-1,
\]
with $a_0=-\infty$ and $a_k=+\infty$ appended so the bins cover $\R$. In code this is
\begin{lstlisting}
interior = frozen.ppf(np.linspace(0, 1, bins + 1)[1:-1])
edges = np.concatenate(([-np.inf], interior, [np.inf]))
\end{lstlisting}
Expected counts are then computed from the CDF directly, rather than assumed to be exactly
$n/k$:
\[
E_i = n\big[F(a_i;\hat\theta) - F(a_{i-1};\hat\theta)\big]
= n \cdot \texttt{np.diff(frozen.cdf(edges))}_i .
\]
Computing $E_i$ this way (rather than hard-coding $n/k$) keeps the statistic internally
consistent when \texttt{cdf} and \texttt{ppf} are not perfectly numerically inverse (e.g.\
for discrete distributions, or distributions using an approximate \texttt{ppf} solver), and
naturally collapses duplicate discrete edges without extra bookkeeping.

\subsubsection{Discrete distributions}
For a distribution with \texttt{pmf} rather than \texttt{pdf}, quantiles can coincide,
producing duplicate edges after \texttt{ppf}. These are merged (\texttt{np.unique}), which
reduces the effective number of bins below the requested $k$ --- an expected consequence of
equiprobable binning on a discrete support, not an error condition. Observed counts on
integer support are computed as
\[
O_i = \#\{x_j \le a_i\} - \#\{x_j \le a_{i-1}\},
\]
implemented via \texttt{np.searchsorted} on the sorted sample, which correctly assigns each
integer-valued observation to the half-open bin $(a_{i-1}, a_i]$.

\subsubsection{Guarding degenerate bins}
The implementation raises if any $E_i \le 0$ or non-finite, which under equiprobable binning
only occurs if the fitted quantile function is degenerate (e.g.\ returns non-increasing
values), signalling an unusable fit rather than attempting to silently patch the bin
structure.

\subsection{Kolmogorov--Smirnov (KS) Statistic}

\subsubsection{Definition}
Let $F_n$ denote the empirical CDF of the sample,
\[
F_n(x) = \frac{1}{n}\sum_{j=1}^n \mathbf{1}\{x_j \le x\},
\]
and let $F(\cdot;\hat\theta)$ be the fitted CDF. The two-sided, one-sample KS statistic is
the sup-norm distance between the two CDFs:
\[
\boxed{\;D = \sup_{x \in \R} \big| F_n(x) - F(x;\hat\theta) \big|\;}
\]
Because $F_n$ is a step function, the supremum is attained at, or immediately to the left
of, one of the observed data points; it suffices to evaluate the discrepancy from both
sides at each observation and take the largest value.

\subsubsection{One-sided components}
Split $D$ into two one-sided statistics that separately capture $F_n$ running above versus
below $F$:
\[
D^+ = \sup_x \big[F_n(x) - F(x)\big], \qquad
D^- = \sup_x \big[F(x) - F_n(x)\big], \qquad
D = \max(D^+, D^-).
\]
For distinct, sorted observations $x_{(1)} < \dots < x_{(n)}$, the classical formulas are
\[
D^+ = \max_{1\le j \le n}\left\{\frac{j}{n} - F(x_{(j)})\right\}, \qquad
D^- = \max_{1\le j \le n}\left\{F(x_{(j)}) - \frac{j-1}{n}\right\}.
\]
Here $j/n$ is the value $F_n$ jumps \emph{to} immediately after $x_{(j)}$ (the ``upper''
empirical value at that point), and $(j-1)/n$ is the value $F_n$ holds \emph{just before}
$x_{(j)}$ (the ``lower'' empirical value). Taking the max/min of $F$ against both bounds at
every observation is what makes the two-sided formulas correct without needing to evaluate
$F_n$ at every real number.

\subsubsection{Handling tied observations}
The formulas above implicitly assume all $x_{(j)}$ are distinct. When observations repeat
(common with rounded, binned, or discrete data), $F_n$ jumps by more than $1/n$ at the
repeated value, and naively using $j/n, (j-1)/n$ indexed by \emph{position} rather than by
\emph{unique value} understates the jump on one side and overstates it on the other,
producing an incorrect $D$. The implementation instead groups by unique value:
\[
v_1 < v_2 < \dots < v_u, \qquad c_i = \#\{x_j = v_i\},
\]
and computes cumulative empirical bounds directly from the tie counts:
\[
F_n(v_i^+) = \frac{1}{n}\sum_{l \le i} c_l, \qquad
F_n(v_i^-) = F_n(v_i^+) - \frac{c_i}{n},
\]
\[
D^+ = \max_i \big[F_n(v_i^+) - F(v_i)\big], \qquad
D^- = \max_i \big[F(v_i) - F_n(v_i^-)\big].
\]
This reduces to the classical formula when all $c_i = 1$, and remains correct when they are
not. In code:
\begin{lstlisting}
values, counts = np.unique(x, return_counts=True)
ecdf_upper = np.cumsum(counts) / n
ecdf_lower = ecdf_upper - counts / n
cdf = frozen.cdf(values)
D = max(np.max(ecdf_upper - cdf), np.max(cdf - ecdf_lower))
\end{lstlisting}

By default, the metrics used for ranking are AIC and BIC. Additional
metrics can be requested by passing a custom \texttt{metrics} list.

\section{Basic Workflow}

The typical workflow uses four steps:

\begin{enumerate}[leftmargin=*]
    \item Load configuration and select candidate distributions.
    \item Construct a \texttt{DistributionFitter} with the data.
    \item Call \texttt{fit()} to estimate parameters and compute metrics.
    \item Review and select the best-fit distribution.
\end{enumerate}

\begin{lstlisting}[style=pythonstyle]
from actsim import DistributionFitter, load_config
from actstats import actuarial as act

# Generate example severity data
sev_data = act.lognormal(0.5, 0.2).rvs(size=10000)

# Load configuration
config = load_config()
distributions = config.distributions["severity"]
metrics = config.metrics

# Construct the fitter
fitter = DistributionFitter(
    sev_data,
    distributions=distributions,
    metrics=metrics,
)

# Fit candidate distributions
fitter.fit()
\end{lstlisting}

\section{Inspecting Results}

After fitting, the best-fit distribution under each metric is stored in
the \texttt{best\_fits} attribute, while the currently selected
distribution is stored in \texttt{selected\_fit}.

\begin{lstlisting}[style=pythonstyle]
print("Best fits:", fitter.best_fits)
print("Selected fit:", fitter.selected_fit)
print("Selected distribution:", fitter.get_selected_dist())
print("Selected parameters:", fitter.get_selected_params())
\end{lstlisting}

\section{Manually Selecting a Distribution}

Users can override the automatic selection by name:

\begin{lstlisting}[style=pythonstyle]
fitter.select_distribution("gamma")
print(fitter.selected_fit["params"])
\end{lstlisting}

This is useful when the analyst has a prior reason to prefer a particular
distribution, or when comparing scenarios.

\section{Sampling and Prediction}

Once a distribution is selected, the fitter can be used as a drop-in
sampler. The \texttt{sample} method returns random draws from the
selected distribution, while \texttt{sample\_mixed} inserts a proportion
of zeros and ones to handle datasets with mass at those points (e.g.\
zero-inflated severity data).

\begin{lstlisting}[style=pythonstyle]
# Draw 10 samples
samples = fitter.sample(size=10)

# Draw 100 samples with 10% zeros and 5% ones
mixed = fitter.sample_mixed(zero_prop=0.1, one_prop=0.05, size=100)

# Evaluate the predicted density at given points
density = fitter.predict([0.5, 1.0, 1.5, 2.0])
\end{lstlisting}

\section{Diagnostic Plots}

The \texttt{plot\_predictions()} method draws a histogram of the data
together with the probability density functions of the fitted candidate
distributions. This plot is the primary visual diagnostic for assessing
fit quality.

\begin{lstlisting}[style=pythonstyle]
fitter.plot_predictions()
\end{lstlisting}

\section{Summary Reports}

Two methods produce numerical summaries:

\begin{itemize}[leftmargin=*]
    \item \texttt{calculate\_statistics()} returns a comparison of the
          mean, standard deviation, and key percentiles between the
          observed data and the predicted density.
    \item \texttt{summary()} returns a DataFrame of all fitted
          distributions, their parameters, and their goodness-of-fit
          metrics.
\end{itemize}

\begin{lstlisting}[style=pythonstyle]
fitter.calculate_statistics()
fitter.summary()
\end{lstlisting}

\section{Practical Notes}

A few practical considerations apply when using \texttt{DistributionFitter}:

\begin{itemize}[leftmargin=*]
    \item Some distributions (e.g.\ Pareto, beta) require positive or
          bounded support. If the input data violate the support, the
          fit for that distribution may fail; \texttt{DistributionFitter}
          catches such failures and continues with the remaining
          distributions.
    \item AIC and BIC are convenient ranking tools but should be
          interpreted alongside diagnostic plots, especially in the
          tails, which often matter most for actuarial applications.
    \item For frequency data, choose distributions appropriate to count
          data (e.g.\ Poisson, negative binomial) rather than continuous
          distributions.
\end{itemize}
